\documentclass[12pt]{article}

\usepackage[a4paper, margin=2.5cm]{geometry}
\usepackage{amsmath}
\usepackage{amssymb}
\usepackage{amsthm}
\usepackage[colorlinks=true, linkcolor=blue, citecolor=blue, urlcolor=blue]{hyperref}
\usepackage{booktabs}
\usepackage{natbib}

\newtheorem{theorem}{Theorem}
\newtheorem{proposition}[theorem]{Proposition}
\newtheorem{corollary}[theorem]{Corollary}
\newtheorem{definition}[theorem]{Definition}
\newtheorem{remark}[theorem]{Remark}
\newtheorem{example}[theorem]{Example}
\newtheorem{lemma}[theorem]{Lemma}

\newcommand{\Fcal}{\mathcal{F}}
\newcommand{\Ecal}{\mathcal{E}}
\newcommand{\Acal}{\mathcal{A}}

\title{Event Structure and the Emergence of Time\\from Phase-Overlap Geometry}
\author{%
  Lee Smart\\[4pt]
  \textit{Institute of Vibrational Field Dynamics}\\[2pt]
  \texttt{contact@vibrationalfielddynamics.org}\\[2pt]
  \texttt{@vfd\_org}%
}
\date{April 2026}

\begin{document}

\maketitle

\noindent\textbf{Status:} Preprint (not peer-reviewed)

\begin{abstract}
Paper~XXII established that the $E_8$ root system contains two conjugate copies of the 600-cell. The present paper investigates the consequence of allowing a relative phase rotation between these two copies.

A relative rotation parameter $\theta$ is introduced as a geometric degree of freedom, not a time coordinate. As $\theta$ varies, pairs of vertices from the two 600-cell copies enter and leave geometric alignment. An event is defined as a vertex pair that becomes admissible for at least one value of $\theta$, and each event is assigned a birth angle: the first value of $\theta$ at which the pair becomes admissible.

Once a reference phase origin $\theta = 0$ is fixed, birth-angle comparison defines a strict order relation on events that is irreflexive and transitive; its non-totality is demonstrated explicitly only in the pentagon toy model of Section~\ref{sec:example}, while in the 600-cell model non-totality is flagged as a separate combinatorial question not settled here. The relation is gauge-dependent in a mild sense: a shift of the reference phase origin leaves the event set fixed (events are static vertex pairs) and cyclically shifts the birth-angle labels assigned to them, so the order is intrinsic to \emph{the geometry together with a chosen phase basepoint}, not to the geometry alone. Conditional on the pentagon-model non-totality and assuming analogous incomparable layers in the 600-cell model, a no-canonical-global-clock statement shows that no unique real-valued total ordering extending $\prec$ can be picked out by its order type alone. A worked example on two pentagons demonstrates the construction explicitly.

The paper does not derive general relativity. It establishes a first step: a causal-type event order derived from discrete closure geometry without assuming spacetime.
\end{abstract}

%==================================================================
\section{Introduction}\label{sec:intro}
%==================================================================

Paper~XXII~\citep{SmartXXII} showed that the $E_8$ root system contains two conjugate copies of the 600-cell (the double cover, 240 roots $\to$ 120 vertex pairs). Papers~IX--X~\citep{SmartIX, SmartX} sketched formal pathways by which continuous geometry and a gravitational analogue might be approached from the constraint manifold, but do not derive full general relativity or a Lorentzian spacetime. The present paper takes a modest next step in the same direction: deriving a causal-type event structure --- and a corresponding emergent ordering interpretable as temporal --- from the phase-overlap geometry of the dual 600-cell system, \emph{once a reference phase origin for the relative rotation parameter $\theta$ has been chosen}.

\subsection*{Core claim}

Time is not assumed as a parameter. An emergent temporal structure arises as the partial ordering of admissible overlap configurations between two phase-rotated copies of the 600-cell constraint geometry.

\subsection*{Scope}

This paper constructs the event structure and proves its ordering properties. It does not derive the Einstein field equations, define a metric, or construct full general relativity. Those are subsequent targets: Paper~XXIV addresses observer frames; Paper~XXV addresses metric emergence; Paper~XXVI addresses curvature; Paper~XXVII addresses discrete graph-Poisson field-equation analogues (not Einstein's equations in the ordinary GR sense). The goal here is to show that, once a reference phase origin has been fixed, a temporal-order-like structure arises from the dual-600-cell geometry.

%==================================================================
\section{The Dual 600-Cell System}\label{sec:dual}
%==================================================================

\subsection{Base geometry}

\begin{definition}[600-cell vertex set]\label{def:H}
Let $H = \{v_1, \ldots, v_{120}\} \subset S^3 \subset \mathbb{R}^4$ be the 120 vertices of the 600-cell on the unit 3-sphere. The adjacency relation is: $v_i \sim v_j$ iff $|v_i - v_j| = 1/\varphi$ where $\varphi = (1+\sqrt{5})/2$.
\end{definition}

The graph Laplacian is $L = 12\,I - A$ where $A$ is the adjacency matrix (degree 12).

\subsection{Dual copies with phase rotation}

\begin{definition}[Dual phase system]\label{def:dual}
The dual 600-cell system consists of:
\begin{itemize}
    \item $H_1 = H$ (the reference copy),
    \item $H_2(\theta) = R_{\theta_0 + \theta}\, H$ where $R_\phi \in SO(4)$ is a fixed one-parameter family of rotations (see remark below) and $\theta_0$ is a fixed generic offset chosen so that at $\theta = 0$ the two copies are in general position (no two vertices $v_i \in H_1$ and $w_j \in H_2(0)$ satisfy $\|v_i - w_j(0)\| = \delta$),
\end{itemize}
with $\theta \in \mathbb{R}$ treated as a lifted real phase parameter, not a time coordinate. Cyclic identification $\theta \sim \theta + 2\pi$ (or whatever the natural period of $R_\phi$ is) is imposed only at the end of any argument that does not depend on the cyclic structure.
\end{definition}

\begin{remark}
The rotation $R_\phi$ acts on $\mathbb{R}^4$ and moves the vertices of $H_2$ on $S^3$. The parameter $\phi$ may be taken along any one-parameter subgroup of $SO(4)$; for definiteness, we use a left-isoclinic rotation $R_\phi(x) = q(\phi) \cdot x$ where $q(\phi) = (\cos\phi/2, \sin\phi/2, 0, 0)$ is a unit quaternion. The generic-position choice of $\theta_0$ excludes measure-zero alignments at the basepoint; this avoids the edge case in which an admissibility interval straddles the phase origin. The qualitative construction below requires only such a one-parameter family of relative configurations and a generic basepoint.
\end{remark}

%==================================================================
\section{Overlap and Admissibility}\label{sec:overlap}
%==================================================================

\subsection{Overlap map}

\begin{definition}[Overlap]\label{def:overlap}
For vertices $v_i \in H_1$ and $w_j \in H_2(\theta)$, define the overlap function
\begin{equation}\label{eq:overlap}
    \Omega(v_i, w_j, \theta) =
    \begin{cases}
        1 & \text{if } \|v_i - w_j(\theta)\| < \delta,\\
        0 & \text{otherwise},
    \end{cases}
\end{equation}
where $\delta > 0$ is a fixed alignment threshold satisfying $\delta < 1/\varphi$.
\end{definition}

\begin{remark}
The condition $\delta < 1/\varphi$ ensures that overlap is a strict proximity condition and does not coincide with adjacency in the 600-cell graph. Thus overlap detects alignment events rather than structural edges.
\end{remark}

\subsection{Admissibility}

\begin{definition}[Admissible pair]\label{def:admissible}
A pair $(v_i, w_j)$ at angle $\theta$ is \emph{admissible} if
\begin{equation}
    \Acal(v_i, w_j, \theta) = \Omega(v_i, w_j, \theta).
\end{equation}
\end{definition}

\begin{remark}
For the purposes of this paper, admissibility is defined purely by geometric alignment. More refined closure constraints may be introduced in later work; here we isolate the event structure generated by overlap alone.
\end{remark}

%==================================================================
\section{Events}\label{sec:events}
%==================================================================

\begin{definition}[Event]\label{def:event}
An \emph{event} is a pair $(v_i, w_j)$ such that there exists at least one $\theta$ with
\[
\Acal(v_i, w_j, \theta) = 1.
\]
\end{definition}

\begin{definition}[Birth angle]\label{def:birth}
Fix the dual system and reference phase basepoint $\theta_0$ as in Definition~\ref{def:dual}. For an event $e = (v_i, w_j)$, define its \emph{birth angle}
\begin{equation}\label{eq:birth}
    \theta_b(e) = \inf\{\theta \geq 0 : \Acal(v_i, w_j, \theta) = 1\},
\end{equation}
where the infimum is taken over the lifted real parameter $\theta \geq 0$ measured from the basepoint. Because the admissibility set $\{\theta : \Acal(v_i, w_j, \theta) = 1\}$ is a finite union of open intervals (strict inequality in Definition~\ref{def:overlap}) and the generic basepoint choice in Definition~\ref{def:dual} places $\theta = 0$ outside all admissibility intervals, the infimum $\theta_b(e)$ is well-defined and equals the left boundary of the first admissibility interval after the basepoint. Events are then understood to be admissible for $\theta$ slightly exceeding $\theta_b(e)$.
\end{definition}

\begin{remark}
Birth angles depend on the chosen reference phase origin: a shift $\theta \mapsto \theta + \Delta$ of the basepoint leaves the event set (the set of vertex pairs that are ever admissible) unchanged, and cyclically shifts every birth-angle label by $-\Delta$ modulo $2\pi$. Events are therefore static geometric objects, but their emergence parameter $\theta_b$ is intrinsic only relative to a chosen phase basepoint.
\end{remark}


Let
\[
\Ecal = \{ (v_i, w_j) : \exists \theta \text{ with } \Acal(v_i, w_j, \theta)=1 \}
\]
be the event set.

%==================================================================
\section{The Ordering Relation}\label{sec:ordering}
%==================================================================

\subsection{Construction}

\begin{definition}[Event ordering]\label{def:ordering}
For events $e_1, e_2 \in \Ecal$, define
\begin{equation}\label{eq:ordering}
    e_1 \prec e_2 \quad\text{iff}\quad \theta_b(e_1) < \theta_b(e_2).
\end{equation}
\end{definition}

\begin{remark}
The parameter $\theta$ is not interpreted as time. The ordering is defined entirely on the event set via the intrinsic emergence parameter $\theta_b$.
\end{remark}

\subsection{Properties}

\begin{proposition}[Strict order; non-totality demonstrated in the pentagon model]\label{prop:partial}
Once a reference phase origin has been fixed, the relation $\prec$ on $\Ecal$ is a strict order: it is irreflexive and transitive. It is non-total \emph{provided} some pair of distinct events shares the same birth angle. Such coincidences occur in the two-pentagon toy model of Section~\ref{sec:example}; whether they occur in the full 600-cell model is a separate combinatorial question that this paper does not settle.
\end{proposition}

\begin{proof}
\textbf{Irreflexivity:} For any $e$, $\theta_b(e) < \theta_b(e)$ is false, so $\neg(e \prec e)$.

\textbf{Transitivity:} If $e_1 \prec e_2$ and $e_2 \prec e_3$, then $\theta_b(e_1) < \theta_b(e_2) < \theta_b(e_3)$, hence $\theta_b(e_1) < \theta_b(e_3)$ and therefore $e_1 \prec e_3$.

\textbf{Conditional non-totality:} If $\theta_b(e_1) = \theta_b(e_2)$ for distinct events $e_1, e_2$, then neither $e_1 \prec e_2$ nor $e_2 \prec e_1$ holds, so the pair is incomparable. The pentagon example (Section~\ref{sec:example}) exhibits such coincidences explicitly. In the 600-cell model the existence of same-birth-angle pairs is a symmetry-driven expectation but not established here.
\end{proof}

\subsection{Interpretation}

\begin{quote}
\textit{The partial order $\prec$ on $\Ecal$ defines an emergent temporal structure.}
\end{quote}

Temporal succession arises from the ordering of admissible configurations. The ordering is not imposed externally; it is determined by the phase-overlap geometry together with a chosen phase origin.

\begin{remark}
The parameter $\theta$ is not interpreted as physical time. It is a coordinate on the space of relative configurations of the dual 600-cell system. The temporal structure arises not from $\theta$ itself, but from the order type induced by the emergence of admissible configurations. In particular, only the ordering of birth angles is used; no metric or duration is derived from $\theta$.
\end{remark}

\begin{remark}
An event represents the minimal admissible interaction channel between the two phase-rotated structures. In this formulation, events are defined at the level of vertex pairs; higher-order interaction structures (edges, faces, or incidence relations) may be incorporated in refined models.
\end{remark}

%==================================================================
\section{No Canonical Global Clock}\label{sec:noclock}
%==================================================================

\begin{theorem}[Linear extensions not unique in the presence of incomparable events]\label{thm:noclock}
Suppose $(\Ecal, \prec)$ contains at least one pair $e_1, e_2$ of incomparable events. Then there exist at least two distinct linear extensions $\leq_1, \leq_2$ of $\prec$, with $e_1 <_1 e_2$ and $e_2 <_2 e_1$. In particular, any order-preserving map $t : \Ecal \to \mathbb{R}$ totally ordering $\Ecal$ is one of an entire family of such maps, none of which is singled out by $(\Ecal, \prec)$ alone.
\end{theorem}

\begin{proof}
Let $e_1, e_2 \in \Ecal$ be incomparable under $\prec$. For each $i \in \{1, 2\}$, let $R_i$ denote the binary relation
\[
  R_1 = \prec \cup \{(e_1, e_2)\}, \qquad R_2 = \prec \cup \{(e_2, e_1)\},
\]
and let $\prec_i$ denote the transitive closure of $R_i$ on $\Ecal$.

\emph{Claim:} each $\prec_i$ is a strict partial order (irreflexive and transitive).

Transitivity is automatic for any transitive closure. For irreflexivity, consider $\prec_1$. Suppose for contradiction that $x \prec_1 x$ for some $x$. Then there is a chain $x = z_0, z_1, \ldots, z_n = x$ in which each consecutive step $(z_j, z_{j+1})$ is either in $\prec$ or equal to $(e_1, e_2)$. If no step uses $(e_1, e_2)$, then all steps are in the strict partial order $\prec$, contradicting irreflexivity of $\prec$. If at least one step uses $(e_1, e_2)$, consider any chain from $e_2$ back to $e_1$ within this cyclic sequence: each segment between consecutive uses of $(e_1, e_2)$ is a $\prec$-chain (after taking transitive closure of $\prec$, which equals $\prec$), so we obtain $e_2 \prec e_1$, contradicting the incomparability of $e_1$ and $e_2$ under $\prec$. Hence $\prec_1$ is irreflexive. The argument for $\prec_2$ is symmetric.

By Szpilrajn's extension theorem~\citep{Szpilrajn1930}, each strict partial order $\prec_i$ on $\Ecal$ extends to a total order $\leq_i$ on $\Ecal$ that also refines $\prec$. By construction $e_1 <_1 e_2$ and $e_2 <_2 e_1$, so $\leq_1 \neq \leq_2$. Any strictly-monotone real-valued map $t_i : (\Ecal, \leq_i) \to \mathbb{R}$ then satisfies $t_1(e_1) < t_1(e_2)$ but $t_2(e_2) < t_2(e_1)$; neither is determined uniquely by $(\Ecal, \prec)$.
\end{proof}

\begin{remark}
This theorem is informational, not metaphysical. It does not say that physical time fails to exist; it says that the discrete order type alone does not single out a preferred linear extension. A physical clock function would require additional structure (e.g.\ a choice of observer worldline or a quantitative alignment-magnitude function on events), not supplied by $(\Ecal, \prec)$ itself.
\end{remark}

\begin{remark}
A total ordering consistent with $\prec$ (a linear extension) may be chosen, but it is not uniquely determined by the event structure. This reflects the absence of a preferred global temporal parameter.
\end{remark}

%==================================================================
\section{Observers}\label{sec:observer}
%==================================================================

\begin{definition}[Observer]\label{def:observer}
An \emph{observer} is a linear extension of $(\Ecal, \prec)$, i.e., a total order $\leq_O$ such that
\[
e_1 \prec e_2 \;\Rightarrow\; e_1 \leq_O e_2.
\]
\end{definition}

\begin{proposition}[Observer multiplicity]\label{prop:observers}
If a set of $k$ events share the same birth angle, then there exist at least $k!$ distinct observers differing in their ordering of these events.
\end{proposition}

\begin{proof}
All events with identical $\theta_b$ are mutually incomparable under $\prec$. Any permutation of these events yields a valid linear extension, giving $k!$ distinct observers.
\end{proof}

\begin{remark}
This provides a combinatorial precursor to observer-dependent simultaneity: all observers agree on ordered events but may differ on the ordering of simultaneous (incomparable) events.
\end{remark}

%==================================================================
\section{Worked Example: Dual Pentagons}\label{sec:example}
%==================================================================

\begin{example}[Two pentagons on $S^1$]\label{ex:pentagon}
Consider the 2D analogue: two regular pentagons on the unit circle $S^1$.

\textbf{Setup.} Pentagon $P_1$ has vertices at angles $\alpha_k = 2\pi k/5$ for $k = 0, \ldots, 4$. Pentagon $P_2(\theta)$ has vertices at $\beta_l(\theta) = 2\pi l/5 + \theta_0 + \theta$, where $\theta_0 = \pi/5$ places $P_2(0)$ midway between adjacent $P_1$-vertices. This makes the minimum $P_1$--$P_2(0)$ distance exactly $\pi/5 > \delta$, so $\theta = 0$ lies strictly outside every admissibility interval (no pair is admissible at the basepoint).

\textbf{Overlap.} For threshold $\delta = \pi/20$ (much smaller than the vertex spacing $2\pi/5$), the overlap $\Omega(\alpha_k, \beta_l, \theta) = 1$ iff $|\alpha_k - \beta_l(\theta)| < \delta$ (mod $2\pi$), equivalently $|2\pi(k-l)/5 - \theta_0 - \theta| < \delta$ (mod $2\pi$).

\textbf{Admissibility.} All overlapping pairs are admissible in this simplified model (no closure constraint beyond proximity).

\textbf{Events.} Pair $(k, l)$ is admissible precisely on the open interval $\theta \in \bigl(2\pi(k-l)/5 - \pi/5 - \delta, \; 2\pi(k-l)/5 - \pi/5 + \delta\bigr)$ (mod $2\pi$). As $\theta$ increases from $0$:
\begin{itemize}
    \item First pairs to become admissible are those with $k - l = 1$: the pair $(k, k{-}1)$ is admissible on $\theta \in (\pi/5 - \delta, \; \pi/5 + \delta)$. All five such pairs $\{(k, k{-}1) : k = 0, \ldots, 4\}$ share admissibility interval and common birth angle $\theta_b = \pi/5 - \delta$. This first layer consists of 5 mutually incomparable simultaneous events.
    \item For $\theta \in (\pi/5 + \delta, \;3\pi/5 - \delta)$, no pair is admissible.
    \item Second layer: pairs with $k - l = 2$, i.e.\ $\{(k, k{-}2) : k = 0, \ldots, 4\}$, each admissible on $\theta \in (3\pi/5 - \delta, \;3\pi/5 + \delta)$ with common birth angle $\theta_b = 3\pi/5 - \delta$.
    \item Further layers at $\theta \approx \pi, 7\pi/5, 9\pi/5$ corresponding to offsets $k - l = 3, 4, 0 \pmod 5$ respectively, each a 5-event layer.
\end{itemize}

\textbf{Ordering.} Layers are totally ordered by their birth angles. Within each layer of 5, no $\prec$-ordering exists (all share the same birth angle). The order on the first two layers is:
\[
    \{(k, k{-}1) : k = 0,\ldots,4\} \;\prec\; \{(k, k{-}2) : k = 0,\ldots,4\} \;\prec\; \cdots
\]

\textbf{No order-type-unique clock.} The 5 events in any one layer are pairwise incomparable, so Theorem~\ref{thm:noclock} applies: any total ordering refining $\prec$ within a layer is an arbitrary choice. Enumerating permutations gives $5! = 120$ distinct linear extensions restricted to each layer.

\textbf{Analogy to the 600-cell (heuristic).} The pentagon has 5 vertices and $\varphi$-geometry (diagonal/side $= \varphi$). The 600-cell is the 4D generalisation with 120 vertices. The pentagon suggests, but does not establish, that analogous same-birth-angle layers may occur in the 600-cell model at symmetry-driven alignment angles; the precise multiplicities are not computed here.
\end{example}

%==================================================================
\section{Relation to Causal Set Theory}\label{sec:causal}
%==================================================================

The event structure $(\Ecal, \prec)$ has formal similarities to causal set theory:

\begin{table}[ht]
\centering
\caption{Comparison with causal set theory.}
\label{tab:causal}
\begin{tabular}{@{}lll@{}}
\toprule
Feature & VFD event structure & Causal set theory \\
\midrule
Events & Admissible vertex pairs & Fundamental elements \\
Ordering & Derived from geometric admissibility (via birth angles) & Assumed as causal relation \\
Origin of set & Derived from 600-cell geometry & Postulated \\
Discreteness & From vertex structure & Postulated \\
Partial order & Proved (Prop.~\ref{prop:partial}) & Assumed \\
Preferred global clock & Not canonically determined (Thm.~\ref{thm:noclock}) & Not fundamental \\
Observer & Linear extension of $\prec$ & Choice of compatible ordering \\
\bottomrule
\end{tabular}
\end{table}

\textbf{Key distinction:} In the VFD framework, the event set and its ordering are \emph{proposed} as derived from the constraint geometry (together with a chosen phase basepoint). In causal set theory, they are postulated as fundamental. The VFD approach proposes a geometric source for an event order that causal set theory takes as primitive.

%==================================================================
\section{Discussion}\label{sec:discussion}
%==================================================================

\subsection{What is established}

\begin{itemize}
    \item The dual 600-cell system with phase rotation is explicitly defined (Definitions~\ref{def:H}--\ref{def:dual}).
    \item Overlap and admissibility are explicitly constructed (Definitions~\ref{def:overlap}--\ref{def:admissible}).
    \item Events are defined as admissible overlap configurations (Definition~\ref{def:event}).
    \item The ordering relation $\prec$ is proved to be a strict order; non-totality is proved explicitly in the pentagon toy model and flagged as conditional in the 600-cell case (Proposition~\ref{prop:partial}).
    \item A no-order-type-unique-global-clock theorem, conditional on the existence of incomparable events, is proved (Theorem~\ref{thm:noclock}).
    \item Observers are defined as linear extensions, with multiplicity proved (Proposition~\ref{prop:observers}).
    \item A worked example demonstrates all constructions explicitly.
\end{itemize}

\subsection{What is not established}

\begin{itemize}
    \item No metric is derived. The partial order gives temporal ordering but not durations.
    \item No Einstein equations. The formal pathways and structural analogies of Papers~IX--X are not upgraded here to field equations; that target is Paper~XXVII.
    \item No Lorentz invariance. The event structure has discrete symmetry ($H_4$), not continuous Lorentz symmetry.
    \item No spatial geometry. The present paper addresses only temporal ordering; spatial structure requires further work.
\end{itemize}

\subsection{Connection to the programme}

The dual 600-cell system of this paper is the same conjugate pair identified in Paper~XXII. The phase rotation $\theta$ is the geometric degree of freedom between Copy~A and Copy~B. The events are the discrete configurations at which the two copies are in admissible alignment. The temporal ordering is the structural consequence of the geometry varying with $\theta$.

This completes the connection:
\begin{itemize}
    \item Papers~I--V: spectral structure (what the constraint geometry produces),
    \item Papers~XII--XXI: quantum dynamics (how it evolves),
    \item Paper~XXII: the spectral bridge (how everything connects),
    \item Paper~XXIII: event structure and emergent time (the first step toward gravity).
\end{itemize}

%==================================================================
\section{Conclusion}\label{sec:conclusion}
%==================================================================

Time is not assumed in the VFD closure framework. An emergent temporal structure is realised as the birth-angle ordering of admissible overlap configurations between two phase-rotated copies of the 600-cell constraint geometry, once a reference phase origin has been chosen. The ordering is irreflexive and transitive (proved); non-totality is proved explicitly in the pentagon toy model and extends to the 600-cell model whenever the generic basepoint still leaves same-birth-angle pairs (an expectation supported by the pentagon analogue but not a theorem of this paper). Conditional on incomparable events, Szpilrajn's extension theorem gives multiple distinct linear extensions, and no one of them is determined by the event structure alone. Different observers correspond to different linear extensions of the partial order, agreeing on ordered events but potentially disagreeing on simultaneous ones.

This is the first step toward gravity within the VFD programme. The subsequent steps --- observer frames (Paper~XXIV), metric emergence (Paper~XXV), curvature (Paper~XXVI), and a discrete graph-Poisson field-equation analogue (Paper~XXVII) --- build on this event structure. The present paper establishes only the order-type-level temporal structure, not the full spacetime geometry.

No external time variable, spacetime manifold, or metric is assumed. The temporal-order-like structure arises from the ordering of admissible configurations within the dual 600-cell geometry together with a chosen phase basepoint. This defines a basepoint-dependent strict event order from the overlap construction; non-totality in the full 600-cell model remains open beyond the pentagon analogue. Subsequent work will address the emergence of metric structure and dynamical laws on top of this event framework.

%==================================================================
\bibliographystyle{plainnat}
\bibliography{references}

\end{document}
